Chúng tôi xác định mười kịch bản theo thông tin bên quan sát muốn suy ra
và phần chuỗi trao đổi (transcript) họ được thấy.

\smallskip\noindent\textbf{Vị trí và đường đã đi.}
Từ một gói, bên quan sát ước lượng GPS hiện tại (S1). Với chuỗi dài hơn,
họ có thể tìm nơi dừng (S2) hoặc khôi phục vị trí ở từng mốc thời gian (S3).
Ở S2, dữ liệu chỉ cho phép dùng tọa độ trung bình trong phút cuối làm đại
diện cho nơi dừng. Khi cho xem toàn phiên, ta còn đo được khả năng suy ra
điểm xuất phát (S9) và điểm kết thúc (S10). Đây là năm tác vụ định vị chính.

\smallskip\noindent\textbf{Liên kết và dự đoán.}
Hai phiên có thuộc về cùng một người hay phương tiện (S4)? Dữ liệu chỉ có
nhãn nhóm tuyến, nên phép thử là xác định hai phiên có cùng nhóm hay không.
Với suy luận tương lai, chúng tôi xét hai quyết định khác nhau trên mạng
đường. S5 dự đoán nhánh kế tiếp trong hai nhánh khả dĩ. S6 chọn một trong
hai đích đến khi đã thấy người dùng rẽ.

\smallskip\noindent\textbf{Nội dung và người đi cùng.}
Ở S7, ta giữ nguyên chuyến đi rồi đổi yêu cầu trên thiết bị để kiểm tra xem
nội dung yêu cầu có làm thay đổi transcript hay không. S8 xét một nguồn
thông tin khác mà nhà cung cấp có thể có. Khi căn hai luồng truy vấn theo
thời gian, quan sát thêm người đi cùng giúp định vị người dùng đến mức nào?
